\subsection{Definition of a locally convex space}

DEFINITION 1. --- A topological vector space is locally convex (real) if there exists a fundamental system of neighbourhoods of 0 that are convex sets.

Such a space is called a locally convex space. Its topology is called a locally convex topology.

The topological vector spaces over \(\mathbf{R}\) which we study in the rest of this book are nearly all locally convex.

If V is a convex neighbourhood of 0 in the locally convex space E, then \(V \cap (-V)\) is a symmetric convex neighbourhood of 0. As the closure of a convex set is convex (II, p.~13, prop. 14) it follows from I, p.~7, prop. 4 that the neighbourhoods of 0 in E which are closed, symmetric and convex, form a fundamental system of neighbourhoods invariant under homotheties of centre 0 and ratio \(\neq 0\) .

PROPOSITION 1. --- Let S be a filter base on a vector space E formed from sets that are absorbent, symmetric and convex. Then the set V of transforms of the sets of S by homotheties of ratio \textgreater{} 0 is a fundamental system of neighbourhoods of 0 for a locally convex topology on E.

Clearly \(\mathfrak{B}\) is a filter base satisfying \((\mathrm{EV}_{\mathrm{I}})\) and \((\mathrm{EV}_{\mathrm{II}})\) of I, p.~7, prop. 4; it also satisfies \((\mathrm{EV}_{\mathrm{III}})\) since if \(\mathbf{V} \in \mathfrak{S}\) then \(\frac{1}{2}\mathbf{V} + \frac{1}{2}\mathbf{V} = \mathbf{V}\) .

Note that if \(\mathcal{T}\) is the locally convex topology on E having \(\mathfrak{V}\) for a fundamental system of neighbourhoods of 0, then the sets \((1/n)\) V, where \(n\) varies in the integers \(>0\) and V varies in \(\mathfrak{S}\) , form a fundamental system of neighbourhoods of 0 for the topology \(\mathcal{T}\) . Then \(\mathcal{T}\) is Hausdorff, if and only if, for every \(x \neq 0\) in E there exists an integer \(n\) and a set \(\mathrm{V} \in \mathfrak{S}\) , such that \(nx \notin \mathrm{V}\) ; if, further, \(\mathfrak{S}\) is enumerable, then the topology \(\mathcal{T}\) is a metrisable locally convex topology. Conversely, it is clear that if \(\mathcal{T}\) is a metrisable locally convex topology, then there exists an enumerable fundamental system of closed symmetric convex neighbourhoods of 0 for \(\mathcal{T}\) .

COROLLARY. --- The topology \(\mathcal{T}\) of a topological vector space E, is defined by a set of semi-norms (II, p.~3) if, and only if, \(\mathcal{T}\) is locally convex.

The condition is necessary since every semi-norm on E is a convex function, and so, for \(\alpha > 0\) , the set of \(x \in E\) for which \(p(x) \leqslant \alpha\) , is convex (II, p.~17, corollary). Conversely if V is a symmetric, closed, convex neighbourhood of 0 in E, the gauge p of V is a semi-norm on E such that V is the set of points x of E satisfying \(p(x) \leqslant 1\) (II, p.~20, prop. 23).

This shows further that a locally convex topology \(\mathcal{T}\) is defined by the set of all semi-norms that are continuous for \(\mathcal{T}\) . Further, if \(\mathcal{T}\) is metrisable, then it is defined by an enumerable set of semi-norms.

From the corollary to prop. 1, all the results of § 1 on topologies defined by sets of semi-norms apply in particular to locally convex topologies over real vector spaces. A locally convex Hausdorff space E has a completion \(\hat{E}\) that is locally convex. A complete, metrisable locally convex space is called a Fréchet space; every Banach space is a Fréchet space.

PROPOSITION 2. --- Let f be a continuous linear form defined over a vector subspace M, of a locally convex space E; then there exists a continuous linear form h that is defined over E and is an extension of f.

From the corollary above and II, p.~7, cor. 2, there exists a continuous semi-norm \(p\) on \(\mathbf{E}\) , such that \(|f(y)| \leqslant p(y)\) for all \(y \in \mathbf{M}\) . By the Hahn-Banach th. (II, p.~23, cor. 1) there exists a linear form \(h\) on \(\mathbf{E}\) that extends \(f\) and is such that \(|h(x)| \leqslant p(x)\) for all \(x \in \mathbf{E}\) , and this implies that \(h\) is continuous (II, p.~6, prop. 5).

Remark. --- If g is a continuous linear mapping of M in the product space \(R^{I}\) , then there exists a continuous linear mapping h of E in \(R^{I}\) that is an extension of g; for writing \(g = (g_{\mathrm{t}})\) , where the \(g_{t}\) are continuous linear forms defined over M, there is an extension \(h_{t}\) of \(g_{t}\) for each \(t \in I\) , such that \(h_{t}\) is a continuous linear form over E. The continuous linear mapping \(h = (h_{\mathrm{t}})\) has the required properties.

Note that if \(\mathbf{F}\) is a locally convex Hausdorff space and \(g\) a continuous linear mapping of \(\mathbf{M}\) in \(\mathbf{F}\) , then there does not necessarily exist a continuous linear mapping of \(\mathbf{E}\) in \(\mathbf{F}\) which is an extension of \(g(\mathrm{IV},\mathrm{p.~55},\mathrm{exerc.~16},c))\) . However there does exist such an extension when \(\mathbf{M}\) is finite dimensional (cf.~cor. 2, below).

COROLLARY 1. --- Let E be a locally convex space. If \(x_0 \in E\) is not in the closure of \(\{0\}\) , then there exists a continuous linear form \(f\) defined over E with \(f(x_0) \neq 0\) .

Apply prop. 2 to the one dimensional vector space M generated by \(x_0\) and to the linear form \(\xi x_0 \mapsto \xi\) defined over M, which, by I, p.~12, prop. 2, is continuous.

COROLLARY 2. --- Let M be a finite dimensional vector subspace of E, a locally convex Hausdorff space. Then there exists a closed vector subspace N of E, which is the topological complement of M in E.

There exists a topological complement to M in E if, and only if, the identity mapping of M on itself can be extended to a continuous linear mapping of E on M, which mapping is then necessarily a continuous projector (GT, III, § 6.2, corollary). Now, this follows from the remark above since M is isomorphic to a space \(R^{n}\) (I, p.~13, th. 2).

PROPOSITION 3. --- In a locally convex space E, the balanced convex envelope of a precompact set is itself a precompact set.

Let A be a precompact set in E. Given V, a balanced convex neighbourhood of 0 in E, there exist finitely many points \(a_i \in \mathbf{A}\) ( \(1 \leqslant i \leqslant n\) ) such that A is contained in S, the union of the neighbourhoods \(a_i + \mathbf{V}\) ( \(1 \leqslant i \leqslant n\) ). Thus C, the balanced convex envelope of A, is contained in T the balanced convex envelope of S; but T is contained in B + V, where B denotes the convex envelope of the finite set of points \(a_i\) , \(-a_i\) ( \(1 \leqslant i \leqslant n\) ). Now B is precompact (II, p.~14, cor. 2); hence there exist finitely many points \(b_k \in \mathbf{B}\) ( \(1 \leqslant k \leqslant m\) ) such that \(\mathbf{B}_k\) is contained in the union of the neighbourhoods \(b_k + \mathbf{V}\) . Then C is contained in the union of the neighbourhoods \(b_k + 2\mathbf{V}\) , and the proposition is proved.

Note that, in an infinite dimensional locally convex Hausdorff space, the convex envelope of a compact set is not necessarily closed (II, p.~74, exerc. 3).

COROLLARY. --- If, in a locally convex Hausdorff space E, a compact set X is contained in a complete convex set (complete in the uniform structure induced by that of E) then the convex closed envelope of X is compact.

For this envelope is a closed subset of a complete space, therefore it is complete, but it is also precompact and Hausdorff.

However in a non complete locally convex Hausdorff space, the convex closed envelope of a compact set need not be compact (II, p.~87, exerc. 2).

